\documentclass[11pt,a4paper]{article}
\usepackage[utf8]{inputenc}
\usepackage[T1]{fontenc}
\usepackage{lmodern}
\usepackage[polish]{babel}
\usepackage{geometry}
\geometry{margin=2cm}
\usepackage{hyperref}
\usepackage{listings}
\usepackage{xcolor}
\usepackage{tcolorbox}

\definecolor{mintdark}{HTML}{2f3542}
\definecolor{techgreen}{HTML}{16a34a}
\definecolor{codebg}{HTML}{f1f2f6}
\definecolor{warningorange}{HTML}{ffa502}

\hypersetup{
    colorlinks=true,
    linkcolor=techgreen,
    urlcolor=techgreen
}

\lstset{
    backgroundcolor=\color{codebg},
    basicstyle=\ttfamily\small\color{mintdark},
    breaklines=true,
    frame=leftline,
    framerule=3pt,
    rulecolor=\color{techgreen},
    framesep=5pt,
    xleftmargin=8pt,
    showstringspaces=false
}

\newtcolorbox{note-box}{
    colback=codebg,
    colframe=techgreen,
    fonttitle=\bfseries,
    title=Uwaga,
    arc=4pt
}

\newtcolorbox{warning-box}{
    colback=codebg,
    colframe=warningorange,
    fonttitle=\bfseries,
    title=Wazna informacja dla Linux Mint,
    arc=4pt
}

\title{\Huge\bfseries\color{mintdark} Poradnik: Instalacja \color{techgreen}ROS 2 \color{mintdark}ze zrodel}
\author{\large Dostosowane dla Linux Mint \& Ubuntu}
\date{}

\begin{document}

\maketitle

\begin{tcolorbox}[colback=codebg, colframe=mintdark]
Czytasz dokumentacje starszej, ale wciaz wspieranej wersji ROS 2. Aby uzyskac informacje o najnowszej wersji, zapoznaj sie z wersja \textbf{\color{techgreen}Lyrical}.
\end{tcolorbox}

\tableofcontents
\newpage

\section{\color{techgreen}1. Wymagania systemowe}
Obecne platformy docelowe oparte na Debianie dla \textbf{Jazzy Jalisco} to:
\begin{itemize}
    \item \textbf{Poziom 1:} Ubuntu Linux - Noble (24.04) 64-bit
    \item \textbf{Poziom 3:} Ubuntu Linux - Jammy (22.04) 64-bit
    \item \textbf{Poziom 3:} Debian Linux - Bookworm (12) 64-bit
\end{itemize}
Zgodnie ze specyfikacja zdefiniowana w \textbf{REP 2000}.

\section{\color{techgreen}2. Konfiguracja systemu}

\subsection{Ustawienia regionalne (Locale)}
Upewnij sie, ze masz ustawienia regionalne obslugujace \textbf{UTF-8}. Jesli pracujesz w minimalnym srodowisku (np. kontenerze Docker), ustawienia moga byc szczatkowe (np. POSIX). Ponizsza konfiguracja to zalecany zestaw testowy.

\begin{lstlisting}[language=bash]
locale

sudo apt update && sudo apt install locales
sudo locale-gen en_US en_US.UTF-8
sudo update-locale LC_ALL=en_US.UTF-8 LANG=en_US.UTF-8
export LANG=en_US.UTF-8

locale
\end{lstlisting}

\subsection{Wlaczenie wymaganych repozytoriow}
Musisz dodac oficjalne repozytorium apt ROS 2 do swojego systemu. Najpierw upewnij sie, ze repozytorium \textbf{Ubuntu Universe} jest wlaczone:

\begin{lstlisting}[language=bash]
sudo apt install software-properties-common
sudo add-apt-repository universe
\end{lstlisting}

Pakiety \texttt{ros-apt-source} dostarczaja klucze oraz konfiguracje dla roznych repozytoriow ROS. Instalacja ponizszego pakietu automatycznie dostosuje system, a jego przyszle aktualizacje beda przebiegac automatycznie.

\begin{lstlisting}[language=bash]
sudo apt update && sudo apt install curl -y
export ROS_APT_SOURCE_VERSION=$(curl -s https://api.github.com/repos/ros-infrastructure/ros-apt-source/releases/latest | grep -F "tag_name" | awk -F'"' '{print $4}')
curl -L -o /tmp/ros2-apt-source.deb "https://github.com/ros-infrastructure/ros-apt-source/releases/download/${ROS_APT_SOURCE_VERSION}/ros2-apt-source_${ROS_APT_SOURCE_VERSION}.$(. /etc/os-release && echo ${UBUNTU_CODENAME:-${VERSION_CODENAME}})_all.deb"
sudo dpkg -i /tmp/ros2-apt-source.deb
\end{lstlisting}

\subsection{Instalacja narzedzi programistycznych}
Uruchom ponizsze polecenie, aby pobrac niezbedne pakiety oraz narzedzia deweloperskie:

\begin{lstlisting}[language=bash]
sudo apt update && sudo apt install -y \
  python3-flake8-blind-except \
  python3-flake8-class-newline \
  python3-flake8-deprecated \
  python3-mypy \
  python3-pip \
  python3-pytest \
  python3-pytest-cov \
  python3-pytest-mock \
  python3-pytest-repeat \
  python3-pytest-rerunfailures \
  python3-pytest-runner \
  python3-pytest-timeout \
  ros-dev-tools
\end{lstlisting}

\section{\color{techgreen}3. Budowanie ROS 2}

\subsection{Pobieranie kodu zrodlowego}
Utworz dedykowana przestrzen robocza (workspace) i sklonuj repozytoria:

\begin{lstlisting}[language=bash]
mkdir -p ~/ros2_jazzy/src
cd ~/ros2_jazzy
vcs import --input https://raw.githubusercontent.com/ros2/ros2/jazzy/ros2.repos src
\end{lstlisting}

\subsection{Instalacja zaleznosci przy uzyciu rosdep}
Pakiety ROS 2 sa budowane na bazie regularnie aktualizowanych systemow. Przed instalacja nowych pakietow upewnij sie, ze Twoj system posiada najnowsze poprawki.

\begin{lstlisting}[language=bash]
sudo apt upgrade
sudo rosdep init
rosdep update
rosdep install --from-paths src --ignore-src -y --skip-keys "fastcdr rti-connext-dds-6.0.1 urdfdom_headers"
\end{lstlisting}

\begin{warning-box}
Jesli uzywasz dystrybucji opartej na Ubuntu, ktora nie identyfikuje sie bezposrednio jako Ubuntu (\textbf{takiej jak Linux Mint}), instalator zglosi blad: \texttt{Unsupported OS [mint]}. W takim przypadku musisz wymusic system docelowy, dodajac flage \textbf{\texttt{--os=ubuntu:noble}} do powyzszego polecenia \texttt{rosdep install}.
\end{warning-box}

\subsection{Dodatkowe implementacje RMW (Opcjonalnie)}
Domyslnym oprogramowaniem posredniczacym (middleware) w ROS 2 jest \textbf{Fast DDS}. Mozesz je jednak podmienic na etapie kompilacji lub w czasie uruchomienia programu.

\subsection{Instalacja colcon mixins}
\begin{lstlisting}[language=bash]
colcon mixin add default https://github.com/colcon/colcon-mixin-repository/raw/master/index.yaml
colcon mixin update default
\end{lstlisting}

\subsection{Kompilacja kodu w przestrzeni roboczej}
Jesli w systemie znajduje sie juz inna instalacja ROS 2 (np. z pakietow binaranych lub deb), upewnij sie, ze uruchamiasz kompilacje w czystym srodowisku. Sprawdz, czy komenda \texttt{printenv | grep -i ROS} zwraca pusty wynik.

\begin{lstlisting}[language=bash]
cd ~/ros2_jazzy/
colcon build --symlink-install --mixin release
\end{lstlisting}

\begin{note-box}
Jesli masz problemy z kompilacja niektorych przykladow i blokuje to caly proces budowania, uzyj flagi \texttt{--packages-skip}. Na przyklad, aby pominac pakiety wymagajace biblioteki OpenCV:
\begin{lstlisting}[language=bash]
colcon build --symlink-install --packages-skip image_tools intra_process_demo
\end{lstlisting}
\end{note-box}

\section{\color{techgreen}4. Konfiguracja srodowiska}
Aby aktywowac zbudowane srodowisko, zaladuj ponizszy plik:

\begin{lstlisting}[language=bash]
. ~/ros2_jazzy/install/local_setup.bash
\end{lstlisting}

\begin{note-box}
Jesli nie uzywasz domyslnej powloki Bash, zastap rozszerzenie \texttt{.bash} odpowiednim dla Twojego shella (np. \texttt{setup.sh}, \texttt{setup.zsh}).
\end{note-box}

\section{\color{techgreen}5. Uruchomienie przykladowych aplikacji}
Otworz \textbf{pierwszy terminal}, zaladuj srodowisko i uruchom wezel nadajacy (C++):
\begin{lstlisting}[language=bash]
. ~/ros2_jazzy/install/local_setup.bash
ros2 run demo_nodes_cpp talker
\end{lstlisting}

Otworz \textbf{drugi terminal}, zaladuj srodowisko i uruchom wezel odbierajacy (Python):
\begin{lstlisting}[language=bash]
. ~/ros2_jazzy/install/local_setup.bash
ros2 run demo_nodes_py listener
\end{lstlisting}

Zobaczysz komunikat \textit{Publishing messages} u nadawcy oraz \textit{I heard those messages} u odbiorcy. Gratulacje! Obie aplikacje komunikuja sie poprawnie.

\section{\color{techgreen}6. Alternatywne kompilatory (Clang)}
Domyslnie uzywany jest kompilator \texttt{gcc}. Jesli chcesz przelaczyc sie na \texttt{Clang}, ustaw odpowiednie zmienne srodowiskowe i wymus rekonfiguracje CMake:

\begin{lstlisting}[language=bash]
sudo apt install clang
export CC=clang
export CXX=clang++
colcon build --cmake-force-configure
\end{lstlisting}

\section{\color{techgreen}7. Odinstalowanie srodowiska}
Dezaktywacja sprowadza sie do otwarcia nowego okna terminala bez ladowania pliku \texttt{setup.bash}. Jesli chcesz calkowicie usunac pliki z dysku i zwolnic miejsce, skasuj katalog roboczy:

\begin{lstlisting}[language=bash]
rm -rf ~/ros2_jazzy
\end{lstlisting}

\end{document}